%%%%%%%%%%%%%%%%%%%%%%%%%%%%%%%%%%%%%%%%%%%%%%%%%
%%%%%%%%%%%% cap: intro %%%%%%%%%%%%%%%%%
%%%%%%%%%%%%%%%%%%%%%%%%%%%%%%%%%%%%%%%%%%%%%%%%%

\chapter{Introduction}\label{cap:intro}


\section{Research Plan}

\subsection{Background and Problem Statement}
Phishing is one of the most utilized vectors in cybersecurity attacks, it has resulted in billions of dollars loses from large scale credential theft, financial fraud and loss of corporate data \cite{Leita2008Phishing,Mokube2007Honeypots}. In multi-stage attacks or advanced persistent threats, the initial compromise of user or company credentials are often obtained through some type of phishing attack.
In order for attackers to continue their operations, they must continue to develop their Phishing strategies to avoid detection and outperform existing cybersecurity companies and other companies providing phishing defense solutions. A large percentage of today's Phishing attacks rely on the use of multiple stages, connected in a chain, to evade current Phishing defenses \cite{Ramachandran2006FastFlux}. Today, Phishing's multi-stage methodology is growing rapidly and it can take many different forms. These multi-stage Phishing campaigns use methods of evasion, redirection, and deception; utilizing Adaptive Payload Delivery, which means the payload or 'baits' are downloaded onto the victim's computer before the attack, and using Real-Time Verification of Victims \cite{Leita2008Phishing}. Spam filters and various other security measures are being developed to help reduce the number of Phishing attacks; however, as these measures are being implemented, Phishing's multi-stage methodology is also evolving. Once Phishing campaigns are detected, it is very difficult for Security Analyst to characterize the technical details of those Phishing campaigns. As a result, for the vast majority of the time users are attacked, Security Analysts do not have a continuous view into the behavior, infrastructure, and post-compromise activities of Phishing attacks. Phishing Honeypots, and other tools being developed to facilitate the detection of Phishing attacks and create Passive Deception, have existed for a limited period of time \cite{Moric2025Honeypots}. Most current Phishing Honeypots or Phishing Detection systems consist of Static and Single-Stage systems, and do not simulate realistic Phishing user interactions \cite{Provos2004ClientHoneypots}. In addition, most current tools for collecting information about Phishing attacks rely heavily on manual collection of data, limiting the number of diverse Phishing behavior collections and reducing the scalability of this data collection method \cite{Bringer2012HoneypotSurvey}.

For many years, honeypots have been utilized as active deception techniques for monitoring malicious activity within controlled environments \cite{Seifert2006Taxonomy}. However most phishing honeypots currently in use do not have the capacity to emulate how a user interacts with a phishing site in a realistic manner or how an attacker will surreptitiously escalate their activities toward a person’s account \cite{Leita2008Phishing}. Further, manual intervention by a human being to an individual’s actions limits the scalability and generalization of the information collected from different phishing behaviours \cite{Bringer2012HoneypotSurvey}.

This thesis addresses these limitations by proposing a multi-layered phishing campaign decoy honeypot incorporating automated malicious user simulation to allow for complete behavioural threat intelligence gathering throughout the entire life cycle of a phishing attack \cite{Moric2025Honeypots}. Ultimately, this research will go beyond traditional detection of phishing attacks, and also generate proactive behavioural threat intelligence.

\subsection{Research Aim and Objectives}

The overarching objective of this research is to design, build, and evaluate a multi-layered phishing honeypot framework that uses enhanced automated malicious user simulation to produce highly trusted behavioural threat intelligence \cite{Moric2025Honeypots}.

The detailed objectives of this research are as follows:

To create a multi-layered phishing decoy infrastructure that simulates actual phishing infrastructures \cite{Leita2008Phishing};

To provide automated behaviour simulation for both attackers and victims enabling the dissemination of the collected information to a large audience with multiple interactions and interface types \cite{Kreibich2005Replay,Alata2006Honeynet};

To collect and catalogue all behavioural indicators relative to a phishing attack, from initial compromise through the final logout \cite{Song2005HoneyStat};

To analyse the patterns of attacker interaction; the evolution of both phishing infrastructure, and evasion tactics utilized \cite{Holz2006AttackPatterns,Freiling2010Evasion};

To assess the effectiveness of the proposed multi-layered phishing honeypot to develop actionable behavioural threat intelligence \cite{Moric2025Honeypots}

To advance the knowledge of information systems and cyber-warfare by the establishment of a new, alternative approach to phishing campaigns that incorporates multiple stages of phishing decoy honeypots as a method to mitigate cyber threats and ultimately secure users' identities \cite{Bringer2012HoneypotSurvey}.

\subsection{Research Questions}

What elements in the honeypot structure design will provide an accurate reflection of multi-phase Phishing Infrastructure as it relates to the real world \cite{Leita2008Phishing,Provos2004ClientHoneypots}? How may Automation of Attacker Simulation allow for increased richness and quantity of behavioral /Tactical Intelligence (data) captured \cite{Kreibich2005Replay,Alata2006Honeynet}? What kind(s) of behavioral indicators may be extracted from Multi-Phase Attacker Interactions \cite{Song2005HoneyStat,Holz2006AttackPatterns}? How well does the proposed framework predict ongoing Phishing techniques, tactics, and methods (TTPs) \cite{Freiling2010Evasion,Bringer2012HoneypotSurvey}?

\subsection{Research Methodology}

This research is based on an experimental /system-oriented methodological approach where Software Development, Controlled Deployment, Automated Interaction (Interaction with a security context through a web browser), and Behavioral Data Analysis will be combined to produce actionable intelligence during the cycle of information collection \cite{Bringer2012HoneypotSurvey}.

System Design

A modular Phishing Honeypot Architecture will be developed and will consist of:

Multi-phase Phishing Web Pages (web sites designed for Phishing /Social Engineering) \cite{Leita2008Phishing}.

A back-end Credential Harvesting simulation, to mimic the use of common credential harvesting websites \cite{Provos2004ClientHoneypots}.

A Redirection and Verification phase to capture all Phishing related activity that occurs while users are redirected from a targeted web site and attempting to verify access to the back-end \cite{Ramachandran2006FastFlux}.

The Data Logging and Telemetry module to gather information about Phishing user activity and behavior \cite{Song2005HoneyStat}.

Automated Attackers will be used within the Phishing Honeypot using:

Headless Browsers, with the goal of providing realistic attacker/user interaction \cite{Provos2004ClientHoneypots};

Scripted models of behavioral simulation that simulate the human nature of input, navigation pattern, and response timing \cite{Kreibich2005Replay,Alata2006Honeynet};

IP Rotation and session diversity will simulate Distributed Attackers \cite{Rajab2007Botnet,Bailey2009Botnets}.

Data Collection

The Phishing Honeypot will log:

HTTP Request/Response and Credential Submission Attempts from Users \cite{Song2005HoneyStat};

Redirection Paths from User to Back-End Credential Harvesters, and Users that made a Response to Verification Requests \cite{Leita2008Phishing};

User-Agent Strings and the Time Characteristics associated with each User-Agent \cite{Holz2006AttackPatterns};

Malware Payloads that are delivered within Phishing Messenger Web Pages \cite{Freiling2010Evasion}. Data Analysis

Analyzed using the following methods:

Behavioral Clustering \cite{Song2005HoneyStat};

Temporal Analysis \cite{Holz2006AttackPatterns};

A Profile of the Comparison of Attack Stages \cite{Holz2006AttackPatterns};

Correlation with Existing Threat Intelligence \cite{Bailey2009Botnets,Rajab2007Botnet}.

\subsection{Tools and Technologies}

A web-based phishing decoy, a system honeypot, automated attacker simulation, and data analysis are all part of the experimental environment of this project. The main honeypot is Cowrie, an SSH/Telnet honeypot that simulates a vulnerable system and captures attacker's activities, including authentication failures, command line inputs, and file transmission.

Flask or Django are the backend technologies used to develop the phishing decoy. The phishing decoy is hosted behind either Nginx or Apache and provides full logging of every request that it receives. Automated attacker simulations and data analysis components will interact with the phishing decoy and Cowrie using a Python-based automated system. This automated system includes HTTP-based attacking via libraries like 'requests' as well as SSH interactions via Paramiko, while selenium or playwright may be used to provide more realistic human-like browser interactions.

All application and system logs are gathered on a central log-collection server that has Filebeat and Logstash configured to capture and send logs to Elasticsearch or OpenSearch for storage, as well as to support visualisation through Kibana. In addition to logs, structured datasets representing behavioural activity may also be stored in PostgreSQL.

The data processing for machine learning is done through the use of pandas and NumPy, while the model building is performed using scikit-learn, PyTorch, or TensorFlow/Keras. The automated forensics component utilizes a Python-based analysis pipeline for reconstructing attack timelines, assigning classifications to incidents, and producing structured incident response reports automatically using Jinja2 and LaTeX/HTML templates. The entire platform has been deployed into a secure virtual environment built with the use of Docker containerization and virtualization on a Linux server.

The technologies and tools presented in this chapter may change during the phases of the project.

\subsection{Ethical and Legal Considerations}

Phishing Decoy Operations will be conducted in a contained experimental environment. No real user data will be collected. No Innocent Victims will be targetted by honeypots. All interactions with the honeypot will be automated simulations or limited within institutional testing environments. The research must conform to your institution's cybersecurity ethics guidelines and to the responsible disclosure guidelines \cite{Bringer2012HoneypotSurvey,Seifert2006Taxonomy}.

The Expected Outcomes are:

A Comprehensive Multi-Stage Phishing Honeypot Platform \cite{Moric2025Honeypots};

A Phishing Interaction Logset of Behavioural Phishing Interactions from the Honeypot \cite{Song2005HoneyStat};

Behavioural Profiles of the Phishing Attacker \cite{Holz2006AttackPatterns};

A Reusable Framework for Future Research on Cyber Deception \cite{Moric2025Honeypots,Bringer2012HoneypotSurvey}.

One or More Publishable Papers from the Empirical Research of the Phishing Honeypot Platform \cite{Moric2025Honeypots}.






\subsection{Work Plan and Timeline}

\begin{figure}[htbp]
    \centering
    \resizebox{\textwidth}{!}{
        \begin{ganttchart}[
            vgrid={*{1}{draw=black!10, dotted}},
            hgrid,
            time slot format=isodate-yearmonth,
            time slot unit=month,
            x unit=1.8cm,
            y unit title=0.8cm,
            y unit chart=0.7cm,
            title label font=\bfseries\footnotesize,
            bar label font=\footnotesize,
            group label font=\bfseries\footnotesize,
            milestone label font=\itshape\footnotesize,
            bar label node/.append style={align=right, text width=5cm, anchor=east, left=5pt},
            group label node/.append style={align=right, text width=5cm, anchor=east, left=5pt},
            milestone label node/.append style={align=right, text width=5cm, anchor=east, left=5pt},
            bar/.style={fill=blue!40, draw=blue!70, rounded corners=2pt},
            bar height=0.5,
            group/.style={fill=black!70, draw=none},
            group height=.2,
            milestone/.style={circle, fill=red, draw=red!50, inner sep=2pt},
        ]{2025-10}{2026-07}

        \gantttitlecalendar{year, month=shortname} \\

        \ganttgroup{Phase 1: Foundation}{2025-10}{2026-01} \\
        \ganttbar{Lit. Review \& Threat Model}{2025-10}{2025-12} \\
        \ganttbar{System Architecture Design}{2025-11}{2026-01} \\

        \ganttgroup{Phase 2: Implementation}{2026-01}{2026-04} \\
        \ganttbar{Honeypot Dev (Web/Backend)}{2026-01}{2026-03} \\
        \ganttbar{Attacker Sim (Bots)}{2026-02}{2026-03} \\
        \ganttmilestone{Integration Complete}{2026-03} \\

        \ganttgroup{Phase 3: Experimentation}{2026-04}{2026-05} \\
        \ganttbar{Deployment \& Collection}{2026-04}{2026-05} \\
        \ganttbar{Data Analysis \& Eval}{2026-05}{2026-05} \\

        \ganttgroup{Phase 4: Completion}{2025-11}{2026-07} \\
        \ganttbar{Dissertation Writing}{2025-11}{2026-06} \\
        \ganttmilestone{Final Defense}{2026-07}

        \ganttlink{elem1}{elem2}
        \ganttlink{elem2}{elem4}
        \ganttlink{elem4}{elem6}
        \ganttlink{elem6}{elem8}
        \ganttlink{elem8}{elem9}

        \end{ganttchart}
    }
    \caption{Project Timeline and Work Plan (2025-2026)}
    \label{fig:gantt_fixed}
\end{figure}